\documentclass[11pt]{letter}
\usepackage[margin=1in]{geometry}
\usepackage{amsmath,amssymb}

\signature{Raeez Lorgat\\
Perimeter Institute for Theoretical Physics\\
31 Caroline Street North\\
Waterloo, ON N2L 2Y5, Canada\\
\texttt{rlorgat@perimeterinstitute.ca}}

\date{\today}

\begin{document}

\begin{letter}{The Editors\\
\emph{Inventiones Mathematicae}}

\opening{Dear Editors,}

I am writing to submit the manuscript ``Modular Koszul Duality and
the Shadow Tower: Algebraicity, Depth Classification, and the
Gravitational Coproduct'' for publication in \emph{Inventiones
Mathematicae}.

The paper extends classical Koszul duality from graded algebras
at a point to chiral algebras on algebraic curves, over the moduli
of stable curves~$\overline{\mathcal{M}}_{g,n}$.  Four results
are proved.  First, the shadow obstruction tower---the sequence of
finite-order projections of the universal Maurer--Cartan
element~$\Theta_{\mathcal{A}}$---is algebraic of degree~$2$ along
each primary line, governed by a shadow metric
$Q(t) = (2\kappa + 3\alpha t)^2 + 2\Delta\, t^2$ whose three seed
invariants $(\kappa, \alpha, S_4)$ determine all higher
coefficients.  Second, the critical discriminant
$\Delta = 8\kappa S_4$ partitions the standard landscape into four
depth classes: Gaussian ($r_{\max} = 2$), Lie ($r_{\max} = 3$),
contact ($r_{\max} = 4$), and mixed ($r_{\max} = \infty$).  Third,
the genus-$g$ free energy satisfies
$F_g = \kappa \cdot \lambda_g^{\mathrm{FP}}$ for uniform-weight
algebras, while multi-weight algebras exhibit a genuine cross-channel
correction ($\delta F_2(\mathcal{W}_3) = (c{+}204)/(16c)$),
resolving the multi-generator universality problem negatively.
Fourth, the gravitational coproduct obtained by BRST reduction is
strictly primitive at all arities.

To my knowledge, this is the first complete modular Koszul duality
theory on algebraic curves, unifying and extending work of
Beilinson--Drinfeld, Francis--Gaitsgory, Loday--Vallette, and
Costello--Gwilliam.

This manuscript has not been published elsewhere and is not under
consideration by any other journal.

\closing{Sincerely,}

\end{letter}
\end{document}
